% PROPOSAL: the author rewrites this section.
\section{Discussion}\label{sec:discussion}

The null is the outcome an efficient market predicts, and it is important to be
precise about which question was asked. LNG cargo tracking is a commercial industry.
Several vendors sell satellite-augmented AIS with proprietary destination inference
to gas trading desks in real time. The signals in this paper are built from public
terrestrial AIS and a free voyage archive, which is a coarser version of what the
marginal trader already has. The test was therefore not whether the physical signal
is informative, but whether public physical data is mispriced in a liquid benchmark
at one- to four-week horizons. A negative answer to that question is close to what
one would expect in advance. What the paper adds is evidence that the negative
result is not an artefact of a defective pipeline, an unlucky specification, or a
search that stopped at the first $t>2$. The pipeline is validated against a
published total, the specification was fixed before the fit, and the search was
corrected for multiplicity and held out.

Three results stand independently of the null. First, the physical nowcasts
converge on the naive mean from above, and for the one-week loadings target this
follows from theory rather than from the particular models tried. The argument
applies to any slow industrial flow series and should temper the expectation that a
richer model will beat a moving average on one. Second, the two construction hazards
of Section~\ref{sec:panel} apply to voyage-derived features in general, and they are
easy to miss even when dates appear to have been handled correctly. The vintage log
that detects them is inexpensive to maintain and, to our knowledge, rarely kept.
Third, the outage rule of Section~\ref{sec:parta} is a usable monitor with a stated
recall.

Several changes could alter the answer. A private information advantage, such as
satellite coverage of the open ocean or a destination model that resolves the
three-quarters of the stock this pipeline cannot, would change the question rather
than the answer and would make the study a test of the vendor's product. Within
public data, the remaining untested approaches are joint rather than marginal fits
(spike-and-slab or partial least squares over the collinear block), a fitted regime
model in place of an observable split, distributed lags of the signals, and the
twenty-six signal keys the panel does not use. Each has a weaker prior than the three
mechanisms already tested, and each faces a larger multiplicity problem against a
residual that has never exceeded \bMaxRsqCov{} of variance. Any of them should be
accompanied by its own pre-registration, its own multiplicity correction, and the
untouched-holdout standard used here.
